%%%%%%%%%%%%%%%%%%%%%%%%%%%%%%%%%%%%
\subsection{Admissibility and Boundary Learning in High Dimensions}
%%%%%%%%%%%%%%%%%%%%%%%%%%%%%%%%%%%%

The admissible set $\Omega$ is the set of states from which a competitive equilibrium can be sustained. As in the RA model, $\Omega$ is unknown a priori and depends on the policy—creating a two-level fixed-point problem.

\subsubsection{Simplifying Assumption: Exogenous Consumption Bounds}

To make the boundary learning tractable, we assume exogenous bounds on consumption:
\begin{equation}
    c^e, c^u \in [c_{\min}, c_{\max}]
\end{equation}
These bounds are set based on economic reasoning (e.g., subsistence consumption below, resource constraints above) and remain fixed throughout training.

With this assumption, the \textbf{endogenous admissible set} reduces to three dimensions:
\begin{equation}
    \Omega_{K,a} = \{(K, a^e, a^u) : \text{feasible equilibrium exists for some } (c^e, c^u) \in [c_{\min}, c_{\max}]^2\}
\end{equation}

\subsubsection{Admissibility Score Components}

The score $\mathcal{A}(s)$ has three components, using the power barrier function from Section 3.2:

\begin{enumerate}
    \item \textbf{Capital Feasibility ($A_K$):}
    \begin{equation}
        A_K = \mathcal{S}(K'; [K_{\min}, K_{\max}], \delta_K)
    \end{equation}
    Ensures next-period capital lies within economically meaningful bounds.
    
    \item \textbf{Asset Feasibility ($A_a$):}
    \begin{equation}
        A_a = \min_{i \in \{e, u\}} \mathcal{S}(a'^i; [0, K'], \delta_a)
    \end{equation}
    Checks that derived assets satisfy borrowing constraints and do not exceed aggregate capital.
    
    \item \textbf{Bond Price Feasibility ($A_Q$):}
    \begin{equation}
        A_Q = \mathcal{S}(Q; [0, \beta], \delta_Q)
    \end{equation}
    Ensures the equilibrium bond price lies in the economically valid range $[0, \beta]$. A price $Q > \beta$ would imply negative real interest rates below the discount rate; $Q < 0$ is economically meaningless.
\end{enumerate}

\textbf{Global Score:}
\begin{equation}
    \mathcal{A}(s) = w_K A_K + w_a A_a + w_Q A_Q
\end{equation}

These component scores provide continuous values in $[0, 1]$ based on proximity to constraint boundaries, enabling smooth gradient signals during training.


%%%%%%%%%%%%%%%%%%%%%%%%%%%%%%%%%%%%
\subsubsection{The Boundary Learning Problem}
%%%%%%%%%%%%%%%%%%%%%%%%%%%%%%%%%%%%

Given a set of sampled states $\{s_i\}_{i=1}^N$ with computed admissibility scores $\{\mathcal{A}(s_i)\}$, we partition into admissible and inadmissible sets based on thresholds. The goal is to learn the boundary of the admissible region $\Omega_{K,a}$ in the three-dimensional space $(K, a^e, a^u)$.

\paragraph{Challenge: High-Dimensional Boundary Representation.}
In the representative agent model, we estimated boundaries via binning in 2D: $B_{\min}(\mu, g), B_{\max}(\mu, g)$. This approach does not scale well to higher dimensions due to the curse of dimensionality—with 20 bins per dimension, 3D would require $20^3 = 8000$ bins, many of which would be sparsely populated.

\paragraph{Solution: $\alpha$-Convex Hull.}
We adopt the \textbf{$\alpha$-shape} (also called $\alpha$-convex hull) to represent the admissible set boundary. This geometric approach:
\begin{itemize}
    \item Directly represents the boundary as a simplicial complex
    \item Handles non-convex regions (unlike the standard convex hull)
    \item Provides efficient membership testing
\end{itemize}

\paragraph{Binary Scoring for Boundary Membership.}
Unlike the component scores $A_K, A_a, A_Q$ which are continuous, the $\alpha$-shape provides a \textbf{binary classification}:
\begin{equation}
    \mathbf{1}_{\Omega}(K', a'^e, a'^u) = \begin{cases}
        1 & \text{if } (K', a'^e, a'^u) \in \mathcal{A}_\alpha \\
        0 & \text{otherwise}
    \end{cases}
\end{equation}
where $\mathcal{A}_\alpha$ is the $\alpha$-shape constructed from admissible points.

Computing a continuous score based on distance to the $\alpha$-shape boundary would be computationally expensive, requiring nearest-boundary-point queries for each sample. The binary indicator is both efficient and sufficient for the sampling procedure: points inside $\mathcal{A}_\alpha$ are candidates for training; points outside are excluded or assigned to $\mathcal{S}_{\text{fail}}$.


%%%%%%%%%%%%%%%%%%%%%%%%%%%%%%%%%%%%
\subsubsection{$\alpha$-Shapes: Definition and Properties}
%%%%%%%%%%%%%%%%%%%%%%%%%%%%%%%%%%%%

The $\alpha$-shape generalizes the convex hull to capture non-convex regions.

\paragraph{Intuition.}
Imagine carving out the point set with a spherical ``ice cream scoop'' of radius $1/\alpha$:
\begin{itemize}
    \item $\alpha \to 0$: The scoop is infinitely large; recovers the \textbf{convex hull}
    \item $\alpha \to \infty$: The scoop is infinitesimally small; recovers the \textbf{point set itself}
    \item Intermediate $\alpha$: Captures concavities, indentations, and potentially holes
\end{itemize}

\paragraph{Formal Definition.}
For a point set $P \subset \mathbb{R}^3$ and parameter $\alpha > 0$:
\begin{enumerate}
    \item Compute the \textbf{Delaunay triangulation} $\mathcal{D}(P)$—a partition of the convex hull into tetrahedra with the empty circumsphere property
    \item For each simplex (tetrahedron, triangle, edge) in $\mathcal{D}(P)$, compute its \textbf{circumradius} $r$
    \item The \textbf{$\alpha$-complex} $\mathcal{A}_\alpha$ consists of all simplices with circumradius $r \leq 1/\alpha$
    \item The \textbf{$\alpha$-shape} is the boundary of the union of simplices in $\mathcal{A}_\alpha$
\end{enumerate}

\paragraph{Membership Testing.}
To test if a point $x = (K, a^e, a^u)$ lies inside the $\alpha$-shape:
\begin{enumerate}
    \item Locate the Delaunay tetrahedron containing $x$ (using spatial search)
    \item Check if this tetrahedron belongs to the $\alpha$-complex
    \item Return \texttt{True} if inside, \texttt{False} otherwise
\end{enumerate}
With appropriate spatial indexing (e.g., AABB trees), membership queries are $O(\log n)$.

\paragraph{Choosing $\alpha$.}
The parameter $\alpha$ controls boundary tightness:
\begin{itemize}
    \item \textbf{Too small} ($\alpha \to 0$): Boundary approaches convex hull, may include infeasible regions
    \item \textbf{Too large}: Boundary is too tight, may exclude feasible regions or create artificial holes
\end{itemize}

\textit{Heuristic:} Set $\alpha$ based on typical point spacing:
\begin{equation}
    \alpha \approx \frac{1}{2 \cdot \text{median}(d_{\text{NN}})}
\end{equation}
where $d_{\text{NN}}$ is the nearest-neighbor distance among admissible points.


%%%%%%%%%%%%%%%%%%%%%%%%%%%%%%%%%%%%
\subsubsection{Integration with the Algorithm}
%%%%%%%%%%%%%%%%%%%%%%%%%%%%%%%%%%%%

\paragraph{Constructing the $\alpha$-Shape from Admissible Points.}

\begin{algorithm}[H]
\caption{$\alpha$-Shape Boundary Construction}
\label{alg:alpha_construction}
\begin{algorithmic}[1]
\Require Scored samples $\{(s_i, \mathcal{A}(s_i))\}$, threshold $\tau_{\text{high}}$, shape parameter $\alpha$
\State \textbf{// Extract admissible points in $(K', a'^e, a'^u)$ space}
\State $\mathcal{P}_{\text{adm}} \gets \{(K'_i, a'^e_i, a'^u_i) : \mathcal{A}(s_i) > \tau_{\text{high}}\}$

\State
\State \textbf{// Compute Delaunay triangulation}
\State $\mathcal{D} \gets \text{Delaunay}(\mathcal{P}_{\text{adm}})$

\State
\State \textbf{// Filter to $\alpha$-complex}
\State $\mathcal{A}_\alpha \gets \{T \in \mathcal{D} : \text{circumradius}(T) \leq 1/\alpha\}$

\State
\State \textbf{// Build spatial index for fast membership queries}
\State $\mathcal{I} \gets \text{BuildSpatialIndex}(\mathcal{A}_\alpha)$

\State \Return $\alpha$-shape $(\mathcal{A}_\alpha, \mathcal{I})$
\end{algorithmic}
\end{algorithm}

\paragraph{Membership Function.}
The learned boundary provides binary classification:
\begin{equation}
    \textsc{IsAdmissible}(K', a'^e, a'^u) = \mathbf{1}\{(K', a'^e, a'^u) \in \mathcal{A}_\alpha\}
\end{equation}

\paragraph{Modified Sampling Procedure.}
During adaptive sampling, candidates are filtered using both the continuous component scores and the binary $\alpha$-shape membership:
\begin{equation}
    \mathcal{S}_{\text{safe}} = \{s : \mathcal{A}(s) > \tau_{\text{high}} \text{ and } \textsc{IsAdmissible}(K'(s), a'^e(s), a'^u(s)) = 1\}
\end{equation}


%%%%%%%%%%%%%%%%%%%%%%%%%%%%%%%%%%%%
\subsubsection{The Two-Level Fixed-Point Structure}
%%%%%%%%%%%%%%%%%%%%%%%%%%%%%%%%%%%%

The $\alpha$-shape boundary learning integrates into the nested fixed-point iteration:

\paragraph{Level 1 (Inner Fixed Point):} For fixed policy $\pi_\theta$:
\begin{enumerate}
    \item Sample candidate states from current domain estimate
    \item Compute forward pass: $s \mapsto (K', a'^e, a'^u, Q, \ldots)$
    \item Compute component scores $\mathcal{A}(s)$
    \item Identify admissible points: $\mathcal{P}_{\text{adm}}^{(n)} = \{(K', a'^e, a'^u) : \mathcal{A}(s) > \tau_{\text{high}}\}$
    \item Update $\alpha$-shape: $\mathcal{A}_\alpha^{(n+1)} \gets \text{AlphaShape}(\mathcal{P}_{\text{adm}}^{(n)}, \alpha)$
    \item Repeat until boundary stabilizes: $\mathcal{A}_\alpha^{(n+1)} \approx \mathcal{A}_\alpha^{(n)}$
\end{enumerate}

\paragraph{Level 2 (Outer Fixed Point):} Alternate between:
\begin{enumerate}
    \item \textbf{Train:} Update $(\pi_\theta, V_\phi)$ on states sampled from current $\mathcal{A}_\alpha$
    \item \textbf{Refine:} Run inner loop to update $\mathcal{A}_\alpha$ based on improved policy
\end{enumerate}

As the policy improves, it may discover feasible equilibria in previously unexplored regions, causing the $\alpha$-shape to expand—the same discovery mechanism as in the representative agent model.


%%%%%%%%%%%%%%%%%%%%%%%%%%%%%%%%%%%%
\subsubsection{Practical Considerations}
%%%%%%%%%%%%%%%%%%%%%%%%%%%%%%%%%%%%

\paragraph{Computational Cost.}
\begin{itemize}
    \item \textbf{Delaunay triangulation:} $O(n^2)$ worst case in 3D, $O(n \log n)$ expected
    \item \textbf{Membership query:} $O(\log n)$ with spatial indexing
    \item \textbf{Reconstruction frequency:} Rebuild $\alpha$-shape every $f_{\text{cache}}$ iterations, not every sample
\end{itemize}

\paragraph{Handling Degeneracies.}
Collinear or coplanar points can cause numerical issues in Delaunay triangulation. Mitigation:
\begin{itemize}
    \item Add small random perturbations to coordinates
    \item Use robust computational geometry libraries
\end{itemize}

\paragraph{Implementation.}
\begin{itemize}
    \item \texttt{scipy.spatial.Delaunay} for triangulation
    \item \texttt{alphashape} package for $\alpha$-shape construction
    \item \texttt{scipy.spatial.Delaunay.find\_simplex} for membership testing
\end{itemize}

\paragraph{Adaptive $\alpha$ Schedule.}
Optionally, start with smaller $\alpha$ (looser boundary) and increase during training:
\begin{equation}
    \alpha^{(k)} = \alpha_{\text{init}} + (\alpha_{\text{final}} - \alpha_{\text{init}}) \cdot \min\left(1, \frac{k - K_w}{K - K_w}\right)
\end{equation}
This allows initial exploration before tightening the feasible region estimate.

